\chapter*{اوپن لاجک پروجیکٹ بارے}
\addcontentsline{toc}{chapter}{اوپن لاجک پروجیکٹ بارے}


\textit{اوپن لاجک متن} رسمی ماورائے منطق تے رسمی طریقیاں دی اک کھلے ماخذ
والی، سانجھے کم نال تیار کیتی درسی کتاب اے۔ ایہ درمیانی درجے توں شروع
ہوندی اے (یعنی رسمی منطق دے تعارفی کورس توں پچھوں)۔ بھاویں ایہ ریاضی
دے ماہر نہ ہون والے پڑھنے والیاں لئی اے (خاص طور تے فلسفے تے کمپیوٹر
سائنس دے طالب علماں لئی)، پر ایہ منطقی سختی نال لکھی گئی اے۔

ہُن شامل کجھ موضوعاں دا بیان ہالے پورا نہیں ہویا، تے کئی حصیاں نوں ہالے
وی وڈی نظرثانی دی لوڑ اے۔ اسی اگے چل کے ہور موضوع شامل کرن لئی متن
نوں ودھاؤن دا ارادہ رکھدے آں۔ اسی متن وچ اصطلاحاں دی فرہنگ، ہور پڑھائی
دی فہرست، تاریخی نوٹ، تصویراں، بہتر وضاحتاں، نتائج دی فلسفے، کمپیوٹر
سائنس تے ریاضی لئی اہمیت دسن والے حصے، تے ہور مشقاں تے مثالاں وی
شامل کرنا چاہندے آں۔ جے تہانوں کوئی غلطی لبھے یا تہاڈے کول کوئی صلاح ہووے،
تاں \href{https://github.com/OpenLogicProject/OpenLogic/wiki/Contributing}{مہربانی کر کے پروجیکٹ دی ٹیم نوں دسو}۔

ایہ پروجیکٹ کھلے ماخذ دے جذبے نال کم کردا اے۔ متن صرف مفت دستیاب ای
نہیں، بلکہ اسی LaTeX دا ماخذ وی کریئیٹو کامنز انتساب لائسنس دے تحت
دیندے آں۔ ایہ لائسنس ہر کسے نوں ساڈا کم ڈاؤن لوڈ کرن، ورتن، بدلّن،
نویں ترتیب دین، ہور صورت وچ لے جان تے اگے ونڈن دا حق دیندا اے، بشرطیکہ
اوہ مناسب انتساب کرے۔ ہور معلومات لئی اوپن لاجک پروجیکٹ دی ویب سائٹ
\href{http://openlogicproject.org/}{openlogicproject.org} ویکھو۔
